\section{Whose Directive Is It}
\label{sec:ch07-whose-directive-is-it}

\index{composition!the composer holds its own definitions|(}%
The exported schema declares \code{@cost} in full.
Chapter~\ref{ch:hot-chocolate-zero} printed the definition and it is still
there, at the bottom of the file with the other directive definitions, saying
what the argument is. Its two descriptions are dropped here as everywhere else
in the book, and nothing else is:

\begin{graphqlsdl}
directive @cost(
  weight: String!
) on
  | SCALAR
  | OBJECT
  | FIELD_DEFINITION
  | ARGUMENT_DEFINITION
  | ENUM
  | INPUT_FIELD_DEFINITION
\end{graphqlsdl}

\index{cost@\code{"@cost}!two definitions of one directive}%
\code{weight: String!}, and the value on every field is the string
\code{"10"}. The composer says that value is not a valid \code{Int!}. It is not
reading the definition in front of it. It has one of its own, and where a
directive is one it knows, its own wins.

\index{listSize@\code{"@listSize}!the arity differs}%
The second error says the same thing in the other direction. The document's
\code{@listSize} takes five arguments, one of which the composer's definition
does not have, so an argument that is legal in the schema the service exported
is an unknown argument by the time it is read. I did not read the composer's
definition in full and do not know its arity; what the error establishes is
which argument is missing from it. Two libraries, one directive name, two
definitions:

% \raggedright in the last column, and so \tabularnewline rather than \\:
% the identifier in the bottom cell is 27 characters with no break point, and
% justified it opens gaps across the two lines above it.
\begin{center}
\begin{tabular}{l p{24mm} p{66mm}}
\toprule
Directive & Hot Chocolate emits & \raggedright the composer expects
  \tabularnewline
\midrule
\code{@cost} & \code{weight: String!} & \raggedright \code{weight: Int!}
  \tabularnewline
\addlinespace
\code{@listSize} & five arguments & \raggedright a definition with no
  \code{slicingArgumentDefaultValue} in it \tabularnewline
\bottomrule
\end{tabular}
\end{center}

\index{cost@\code{"@cost}!is not a federation directive}%
Neither directive belongs to federation, which is why neither vendor is obliged
to follow the other. Both come from the cost-analysis work that a client-facing
GraphQL server does to decide how much a query may ask for, and
chapter~\ref{ch:federation-model} already put them in that category: documented
by Apollo alongside the federation directives, implemented by Hot Chocolate in
its core rather than in the federation package, and about how much work a
client may ask for rather than about a seam. Two implementations of an idea
that is still moving have drifted apart, and the drift is invisible until
something reads both.

\index{composition!which side is right was not established}%
I am not going to tell you which of them is right, because I did not establish
it. The disagreement is the fact, and the direction of it is the fact: Hot
Chocolate says string, the composer says integer, and your subgraph is the
place the two meet.

\subsection{One Line}

\index{ModifyCostOptions@\code{ModifyCostOptions}!turns the defaults off}%
The weight's type is the library's and I found nothing that changes it, which
is fine, because the directives were never configured in the first place. Stop
emitting them. That is one call on the builder:

\begin{minted}[highlightlines={35}]{csharp}
using ChilliCream.Nitro.App;
using HotChocolate.AspNetCore;
using Microsoft.EntityFrameworkCore;
using Sessions;

var builder = WebApplication.CreateBuilder(args);

// EF Core reports every command through the logging pipeline as well,
// at Information and with a duration attached. StatementLog below is
// this service's instrument, so silence the other one rather than have
// every statement printed twice.
builder.Logging.AddFilter(
    "Microsoft.EntityFrameworkCore.Database.Command",
    LogLevel.Warning);

builder.Services.AddDbContextFactory<ConferenceContext>(options =>
{
    options.UseSqlite("Data Source=conference.db");

    if (builder.Environment.IsDevelopment())
    {
        options.AddInterceptors(new StatementLog());
    }
});

builder
    .AddGraphQL()
    .AddSessionsTypes()
    .RegisterDbContextFactory<ConferenceContext>()
    .AddQueryContext()
    .AddPagingArguments()
    .AddGlobalObjectIdentification()
    .AddMutationConventions()
    .AddApolloFederation()
    .ModifyCostOptions(options => options.ApplyCostDefaults = false);

var app = builder.Build();

using (var scope = app.Services.CreateScope())
{
    var factory = scope.ServiceProvider
        .GetRequiredService<IDbContextFactory<ConferenceContext>>();
    using var db = factory.CreateDbContext();
    db.Database.EnsureCreated();
}

app.MapGraphQL()
    .WithOptions(options =>
    {
        // Serve the Nitro build that shipped in the package. The
        // default, ServeMode.Latest, downloads the current one
        // from ChilliCream's CDN instead.
        options.Tool.ServeMode = ServeMode.Embedded;
    });

app.RunWithGraphQLCommands(args);
\end{minted}

\index{ApplyCostDefaults@\code{ApplyCostDefaults}!removes both directives}%
\code{ApplyCostDefaults} removes both directives, uses and definitions
together, and the exported schema goes from 201 lines to 165.

\index{ApplyCostDefaults@\code{ApplyCostDefaults}!the option beside it is not the fix}%
\code{ApplySlicingArgumentDefaultValue} is the option next to it, and it looks
like the one for the \code{@listSize} half. It is not. Setting it false drops a
single argument from a directive the composer would reject for another reason
anyway, so the graph still does not compose, and once the defaults are off it
changes nothing at all. One option, one line.

\index{schema export@\code{schema export}!has no filter for this}%
The other place you might look for a fix is the export, and it is not there.
\code{schema export} takes an output path, a schema name and a flag about
semantic nullability, and nothing that drops a directive on the way out. That
is the right design. The exported document is what the service serves from
\code{\_service}, and a switch that made the two differ would be a worse
problem than this one.

\subsection{What That Costs, and What It Does Not}

\index{listSize@\code{"@listSize}!its numbers were chapter 3's}%
Chapter~\ref{ch:data-without-n-plus-1} set a default page of two and a maximum
of ten, and said that the exported file carries those numbers inside
\code{@listSize}. It no longer does, because the directive is gone. Worth being
precise about what that means, because it sounds like a chapter losing a
feature.

\index{paging!the limits still hold}%
The limits still hold. Ask for the schedule with no arguments and two sessions
come back with a next page waiting. Ask for eleven and the request is refused:

\begin{minted}{text}
The maximum allowed items per page were exceeded.
\end{minted}

\index{cost analysis!the directive was a description, not a control}%
Ask for ten, which is the maximum, and all four sessions arrive. So
\code{[UsePaging]} governs paging exactly as it did, and the directive that
disappeared was never what enforced it. It was the cost analyzer's description
of a limit that lives somewhere else. Losing the description costs the schema a
statement about itself and costs a running service nothing.

\index{cost analysis!what turning the defaults off gives up}%
What you do give up is the default weights, which is a real thing to give up if
you were relying on them. A schema with no \code{@cost} anywhere is a schema
where every field costs whatever the analyzer's fallback is rather than what
the library guessed, and a service that wants complexity limits under a
federated router has to say what it wants rather than accept a default.
Chapter~\ref{ch:operations} is where that becomes a decision rather than a
side effect. For now the trade is: a graph that composes, in exchange for
weights nobody chose.
\index{composition!the composer holds its own definitions|)}%
